\section{Related Work}
\label{sec:related}

Concord draws on four lines of prior work: materialized views and result caching, classical multi-query optimization, adaptive physical tuning, and recent LLM-for-database systems. We discuss each in turn, emphasizing the design assumption each makes about how reuse opportunities are discovered.

\subsection{Materialized Views and Result Caching}

Materialized views~\cite{halevy2001mv, goldstein2001mv} and result recyclers~\cite{ivanova2010recycler,nagel2013recycling} maintain precomputed results that later queries can be rewritten against. The core technical challenges have been view-matching---determining whether a stored result can answer a query~\cite{goldstein2001mv}---and maintenance under updates~\cite{gupta1995mvmaint}. MonetDB's recycler~\cite{ivanova2010recycler} and related systems~\cite{phan2008predictableperf} demonstrate that operator-level intermediate caching can yield substantial speedups on workloads with repeated subexpressions.

These systems are designed for workloads in which reuse opportunities must be discovered from observed history. Their decisions are reactive: a result is recycled when its reuse likelihood, estimated from workload statistics, justifies the storage and maintenance cost. Concord inverts this: the agent \emph{declares} which result will be reused, so the materialization decision is made at first execution, without requiring statistical evidence of recurrence. This is not a claim that declared intent is universally better---it is a claim that when the client can be trusted to declare faithfully, a richer optimization space becomes legal. Section~\ref{sec:protocol} makes the legality argument precise.

\subsection{Classical Multi-Query Optimization}

Multi-query optimization (MQO)~\cite{sellis1988mqo,roy2000mqo,dalvi2001pipelined} identifies common subexpressions across a known batch of queries and generates a shared plan that evaluates them together. MQO has been studied both as a static batch problem~\cite{roy2000mqo} and in pipelined~\cite{dalvi2001pipelined} and online~\cite{silva2012onlinemqo} settings. Recent work~\cite{mqo2020revisit} has revisited MQO in the context of modern analytical workloads.

MQO classically assumes a \emph{static} batch of queries known in advance. Concord's M2 mechanism is closer in spirit to pipelined and online MQO: the agent declares a $k$-query lookahead window ahead of execution, and the harness co-optimizes within that window. What distinguishes Concord is not the shared-plan construction itself, which builds on established MQO techniques, but the signal used to identify the batch: an explicit lookahead declaration from the agent, rather than a scheduler buffering concurrent queries. This signal carries additional structure---the agent's declared \emph{diff clauses} between lookahead queries identify exactly which predicates and group-bys vary---that a scheduler observing only the query text cannot recover without further analysis.

\subsection{Adaptive Physical Tuning}

Database cracking~\cite{idreos2007cracking} and its descendants~\cite{idreos2011crackingupdates,petraki2015holistic} adapt physical structure (sort orders, indexes) from query predicates observed at runtime. Soft indexes~\cite{graefe2010softindex} and adaptive merging~\cite{graefe2010adaptivemerge} generalize the idea. These systems, like materialized view work, infer physical structure from observed queries.

Concord's M3 mechanism adapts physical properties of the session workspace---sort order, partial aggregates, session-level GUCs---from the \emph{declared} access pattern in the lookahead window, not from observation. The distinction matters empirically: our index negative result (Section~\ref{sec:eval}) shows that traditional indexes do not help the agent's full-scan analytical workload, because the agent rarely issues point lookups on the workspace. What does help is physical structure aligned with the agent's declared aggregations---partial aggregates, sort orders matching declared \texttt{GROUP BY} keys---which observation-based systems cannot synthesize without seeing the queries first.

\subsection{LLM Agents and Databases}

Recent work has explored LLM agents as database clients in several roles. Text-to-SQL systems~\cite{spider2018,bird2023,spider2_2024} treat the agent as a natural-language-to-SQL translator benchmarked on single-shot translation accuracy. Benchmarks for multi-step analytical agents~\cite{insightbench2024,dabbench2024} evaluate longer-horizon task completion but do not consider database-side protocol adaptation. Systems such as DB-GPT~\cite{dbgpt2023}\TODO{verify cite} and related frameworks provide the agent with tool interfaces over the database but leave the underlying query execution unchanged.

Closely related is the line of work on Automated Data Exploration by Milo and collaborators~\cite{linx2025,atena2020,atenapro2023}, which uses constrained reinforcement learning and LLM-derived specifications to \emph{generate} tree-structured exploration sessions for human analysts. LINX's generated sessions exhibit the same base-and-variants decomposition we observe empirically in agent sessions (Section~\ref{sec:motivation}), providing external corroboration of the pattern, but the optimization target is different: LINX generates sessions offline for later human consumption, while Concord optimizes the execution cost of an agent's sessions in the interactive loop.

Earlier work on next-step recommendation in interactive data analysis~\cite{react2018} infers likely next queries from session logs and suggests them to human users. Concord inverts the signal direction: the LLM agent emits prospective declarations of its own intent, and the database uses those declarations to optimize execution.

Concurrent with our work, several efforts have begun to ask whether databases should adapt to LLM-generated workloads \TODO{cite recent agentic-DB positioning papers}. To our knowledge, Concord is the first system to treat the agent's structured intent---declared reuse, declared lookahead---as the primary optimization signal, and to characterize empirically both the fidelity of agent-emitted declarations and the session-level speedup they enable.

\subsection{Summary of Positioning}

Prior cooperative-database mechanisms treat the client as a black box that emits queries. They infer what to cache, share, or index from observed workload history. Concord treats the LLM agent as a \emph{cooperative} client that can declare intent prospectively, and designs a minimal protocol to exploit that cooperation. The three mechanisms of Section~\ref{sec:protocol} each identify a specific form of intent the agent can declare and a specific optimization that becomes legal only under that declaration.
